\section{Critical points and strict pairing}\label{sec:pairing}

For \(P\)-almost every map and area-almost every target, our next goal is to pair
each saddle of the potential with a higher non-pole local maximum.
Compact superlevels and nondegeneracy will supply finite critical-point
counts above positive levels.

\subsection{The critical map}

Define the smooth real map and its normalized signed Jacobian by
\begin{equation}\label{eq:G}
 G_F(z)=F(z)-\frac{iyF'(z)}k,\qquad
 J_F(z)=\frac{\det DG_F(z)}{|F'(z)|^2}.
\end{equation}
Its real derivative columns are
\((F'/k)(k+q_F)\) and \((iF'/k)(k-1+q_F)\).
Consequently
\begin{equation}\label{eq:Jacobian}
 k^2J_F=k(k-1)+(2k-1)a_F+|q_F|^2,\qquad
 k^2\E J_F(i)=\frac{(4\beta-1)(\beta+2)}9.
\end{equation}
The second identity uses \eqref{eq:moments} and \(k=(\beta+2)/3\).

Off the possible pole, write \(\ell=\log V_\xi^F\) and
\(r_0=F'/(F-\xi)\). Then
\[
 \Delta\ell=-k/y^2,\qquad
 (\ell_x,\ell_y)=(-\re r_0,\ k/y+\im r_0).
\]
Critical points are therefore exactly the preimages \(G_F(z)=\xi\).
A pole is not such a preimage. At a critical point,
\begin{equation}\label{eq:critical-value}
 V_\xi^F(z)=k y^{k-1}|Q_F(z)|,\qquad
 r_0'=\frac{k(k+q_F)}{y^2}.
\end{equation}
Since \(3(k-1)=\beta-1\), cubing the critical value produces
\(k^3y^{\beta-1}|Q_F(z)|^3\). This is the weight, up to the constant
\(k^3\), that the target change of variables in
Section~\ref{sec:transport} will assign to each critical point.
Direct differentiation yields
\begin{equation}\label{eq:hessian}
 \operatorname{Hess}\ell
 =\frac{k}{y^2}
 \begin{pmatrix}-(k+a_F)&b_F\\b_F&k-1+a_F\end{pmatrix},
 \qquad
 \det\operatorname{Hess}\ell=-\frac{k^4}{y^4}J_F.
\end{equation}
Its trace is negative. Thus \(J_F>0\) gives a nondegenerate saddle,
and \(J_F<0\) a nondegenerate local maximum. We call the latter
a \emph{finite maximum}, excluding a pole.

\subsection{Good targets}

Let \(\mathcal G\subset\KK\times\C\) be the set of pairs
\((F,\xi)\) for which all positive closed superlevels
of \(V_\xi^F\) are compact, \(\xi\) is a regular value of \(G_F\),
and distinct preimages under \(G_F\) have distinct \(V_\xi^F\)-values.
We call these pairs \emph{good}.

\begin{lemma}[Generic targets and Borel goodness]\label{lem:good}
Goodness is Borel in \((F,\xi)\). For \(P\)-almost every \(F\),
area-almost every \(\xi\) is good.
\end{lemma}
\begin{proof}
The superlevel assertion is Proposition~\ref{prop:compact-levels}.
For any fixed \(F\), singular values of \(G_F\) have area zero.
This is the plane case of Sard's theorem \cite[Theorem 4.1]{Sard}.
Here is the elementary estimate needed for this smooth plane map.
On a compact source box, cover singular points by grid squares
of side \(h\). Linearize at one singular point of each square.
Bounded second derivatives put its image in an \(O(h^2)\)-neighborhood
of a segment of length \(O(h)\). Its area is \(O(h^3)\).
There are \(O(h^{-2})\) squares, giving total area tending to zero.
Exhaust the source by compact boxes.

Cover the nonsingular locus by countably many inverse-function
charts. On an overlap where two branches \(z_1(\xi),z_2(\xi)\)
are distinct, the target differential of their critical log values is
\[
 d_\xi\bigl[\log V_\xi^F(z_j(\xi))\bigr](v)
 =\re\frac{v}{F(z_j(\xi))-\xi}.
\]
The source derivative vanishes by criticality.
The two linear functionals differ by injectivity of \(F\).
Hence the difference of log values has nonzero differential, and
its zero set is locally a smooth curve. Countably many chart pairs
give an area-null tie set.

For measurability, use
\(C_m=\{|x|\le m,\ 1/m\le y\le m\}\).
For every integer \(n\ge1\), require the open set
\(\{V_\xi^F>1/n\}=\{|F(z)-\xi|<ny^k\}\) to lie in some \(C_m\).
Failure of a given containment is detected at rational-coordinate
points outside the closed box. This is a Borel test.
It is equivalent to compactness of all closed positive superlevels:
for a closed level choose a strictly smaller reciprocal-integer
threshold and use relative closedness.

Failures of regularity or absence of ties are witnessed by one or
two points in a common \(C_m\). For a tied pair impose separation
at least \(1/n\), and test the values by the continuous expression
in \eqref{eq:critical-value}. Equality conditions and this separation
condition have closed parameter projections because the source
boxes are compact. Countably many \(m,n\) prove the other Borel tests.
\end{proof}

If \(t\in\HH\), \(Y=\im t\), and
\(\eta=(\xi-F(t))/(YF'(t))\), direct calculation gives
\begin{align}
 G_{T_tF}(w)&=\frac{G_F(t\cdot w)-F(t)}{YF'(t)},&
 J_{T_tF}(w)&=J_F(t\cdot w),\label{eq:G-action}\\
 V_\eta^{T_tF}(w)&=Y^{1-k}|F'(t)|\,V_\xi^F(t\cdot w).
 &&\label{eq:V-action}
\end{align}
Thus goodness corresponds in both directions under rebasing.

\subsection{The superlevel count}

\begin{lemma}[The superlevel index count]\label{lem:count}
For a good pair and \(h>0\),
\begin{equation}\label{eq:count}
 \#\{z:\ z\text{ is a finite maximum},\ V_\xi^F(z)\ge h\}
 \ \ge\
 \#\{z:\ z\text{ is a saddle},\ V_\xi^F(z)\ge h\}.
\end{equation}
Both counts are finite.
\end{lemma}
\begin{proof}
The fiber \(G_F^{-1}(\xi)\) is closed and discrete, and the closed
positive superlevel is compact. This proves finiteness. There can
be no accumulation at a pole either, because a pole is not on the fiber.

First count in \(U=\{V_\xi^F>\lambda\}\) for a positive regular
level \(\lambda\). We will show that each component is bounded by
one simple closed curve and has total gradient index one, giving
\begin{equation}\label{eq:index-count}
 \#\text{components}
 =\#\text{finite maxima}-\#\text{saddles}+\#\text{poles}
 \quad\text{in }U.
\end{equation}
The boundary of \(U\) is the compact smooth one-manifold
\(\{V_\xi^F=\lambda\}\), away from any pole.
It is a finite disjoint union of smooth simple closed curves.
For detail, the components are open in this compact manifold, so
there are finitely many. Each component is one orbit of the smooth
unit tangent field obtained by rotating \(\nabla\ell/|\nabla\ell|\).
The orbits are complete and open. A nonperiodic single orbit would
be a bijective local homeomorphism from \(\R\) to the compact
component, which is impossible. Its least period gives a simple
closed curve.

The bounded Jordan interior of each curve has closure in \(\HH\):
the region below the curve's minimum height lies in the unbounded
complement. Everywhere inside, \(V_\xi^F>\lambda\). Otherwise
\(\ell\) would have an interior minimum, contrary to
\(\Delta\ell=-k/y^2<0\). If a pole is present, remove a sufficiently
small neighborhood where \(\ell\) is large before taking this minimum.
Thus curves cannot be nested, and their interiors are exactly the
components of \(U\). Indeed each component meets the level boundary,
has its superlevel side locally inside one of the curves, and cannot
leave that interior without crossing the level.

On each positively oriented curve, \(\nabla\ell\) points inward
and has winding one. We recall the tangent-turning argument
of Hopf \cite[Section 2]{Hopf}.
Parametrize a simple regular curve by \(\gamma:[0,1]\to\C\),
with matching endpoint derivatives, starting at a leftmost point.
The positive tangent initially points down. The unit secant direction
from \(\gamma(s)\) to \(\gamma(t)\), \(0\le s\le t\le1\),
extends continuously to the diagonal by the tangent, and to
\((0,1)\) by the negative initial tangent.
Lift its angle continuously on this triangle.
On \((0,t)\) it lies in the closed right semicircle and changes
angle by \(\pi\). On \((s,1)\) it is the negative of the direction
at \((0,s)\), giving another \(\pi\).
The diagonal therefore records tangent winding one, unchanged for
the inward normal.

Remove disjoint small disks about the finitely many critical points
and possible pole in each component. Stokes applied to the closed
angle form
\[
 \frac{X_1\,dX_2-X_2\,dX_1}{|X|^2},\qquad X=\nabla\ell,
\]
shows that the sum of interior indices is one.
At a nondegenerate critical point the index is the sign of the
Hessian determinant, hence \(+1\) for a maximum and \(-1\) for
a saddle. At the pole \(z_0\),
\(\ell=-\log|z-z_0|+\text{a smooth function}\), giving index \(+1\).
Summing over components proves \eqref{eq:index-count}.
There is at most one pole; if present it belongs to a component.
Thus the finite-maximum count is at least the saddle count.
This also covers the empty superlevel.

Finally choose a regular \(\lambda<h\) with no critical value
in \([\lambda,h)\). Such a choice is possible by finiteness
above \(h/2\). The counts in \(U\) are then the counts in
\eqref{eq:count}.
\end{proof}

\subsection{Measurable rank pairing}

The count gives an intrinsic assignment: match the two decreasing
lists by rank. Measurability requires a separate argument, because
we will integrate the assignment over maps and source points. General
countable-section uniformization is supplied by the Lusin--Novikov
theory; see \cite[4F.6 and 4F.17]{Moschovakis}. In the present
setting, local inverse branches give an elementary enumeration, while
intrinsic ranks ensure equivariance.

\begin{proposition}[Strict measurable rank pairing]\label{prop:pair}
On the Borel domain
\[
 \mathcal D=\{(F,z):J_F(z)>0,\ (F,G_F(z))\in\mathcal G\}
\]
there is a Borel assignment \(R_F(z)\) to a finite maximum over
the same target, injective on each target fiber, such that
\begin{equation}\label{eq:strict-value}
 V_{G_F(z)}^F(R_F(z))>V_{G_F(z)}^F(z).
\end{equation}
It is equivariant:
\[
 R_F(t\cdot w)=t\cdot R_{T_tF}(w)
\]
on the corresponding domains.
\end{proposition}
\begin{proof}
On each good fiber order saddles by decreasing critical value,
and separately order finite maxima the same way.
Each point has finitely many predecessors. A nonempty list has a
first point: choose any one point and inspect the finite nonempty
set at or above its value. Repeating gives integer ranks for all points.
Lemma~\ref{lem:count} supplies at least as many maxima as saddles
at or above each saddle's value. Pair the saddle of rank \(j\)
with the maximum of rank \(j\).
No ties are allowed, so the paired value is strictly larger,
as illustrated in Figure~\ref{fig:ranks}.
Different ranks give different partners. Equation~\eqref{eq:V-action}
preserves ranks, proving equivariance.

We prove the measurability assertion directly using the nonsingularity
of each root.
The nonsingular solution graph \(G_F(z)=\xi,\ J_F(z)\ne0\)
has countably many continuous local root branches \(b_j(F,\xi)\),
defined on open parameter sets.
To obtain a branch near \((F_0,\xi_0,z_0)\), set
\(M=DG_{F_0}(z_0)\). On a sufficiently small closed source ball,
\[
 w\longmapsto w-M^{-1}(G_F(w)-\xi)
\]
has Lipschitz constant at most \(1/2\) for all nearby parameters.
Shrinking the parameter neighborhood makes the displacement of
the center less than half the radius. The map is a contraction
of the ball into itself, with a unique interior nonsingular root.
Compactness and uniqueness imply continuous dependence on the
parameters. These branches form relatively open graph patches;
second countability of \(\KK\times\C\times\HH\) gives a countable cover.

At each parameter activate branch \(j\) only when it is defined
and its value differs from every earlier defined branch value.
This Borel test enumerates each root once.
The signs of \(J\), critical values \eqref{eq:critical-value},
and comparisons of such values are Borel. On good fibers,
counting active predecessors of the same sign gives Borel integer
ranks. Selecting the unique active negative-sign branch of the
same rank as a given active positive-sign branch is therefore Borel.
Substitute \(\xi=G_F(z)\) to obtain the claimed assignment.
\end{proof}

\begin{figure}[tbp]
\centering
\begin{tikzpicture}[x=1cm,y=.65cm,
 point/.style={circle,fill,inner sep=1.5pt},
 >={Stealth[length=5pt]}]
 \node at (0,5.6) {Saddles};
 \node at (3.5,5.6) {Finite maxima};
 \foreach \j/\sy/\my in {1/4.1/4.8,2/2.3/3.7,3/.7/2.9}{
   \node[point,label=left:{$s_{\j}$}] (s\j) at (0,\sy) {};
   \node[point,label=right:{$m_{\j}$}] (m\j) at (3.5,\my) {};
   \draw[->,semithick] (s\j) -- (m\j);
 }
 \node[point,label=right:{$m_4$}] at (3.5,1.6) {};
 \draw[->] (5.6,0) -- (5.6,5.1) node[above] {$V_\xi^F$};
\end{tikzpicture}
\caption{Schematic rank ordering on one good fiber \(G_F^{-1}(\xi)\).
Both lists are ordered by decreasing critical value, and equal ranks
are paired. Distinct critical values put each paired maximum strictly
above its saddle. Heights represent values rather than spatial
positions; the displayed counts and gaps are illustrative.}
\label{fig:ranks}
\end{figure}

